\emph{Setting.} Every statement is pointwise algebra except the elimination step and (iv). At a point use the displacement gauge in which the coframe $\hb^{-1}$ is the identity on an orthonormal frame $w_{0}, \dots, w_{3}$ of the derived metric ($Q(w_{0}) = 1$, $Q(w_{i}) = -1$); all quantities are covariant, so this loses nothing. With $\Theta(a) = -\nInf \wedge a + \nO \wedge S(a) + \beta(a)\,\nInf \wedge \nO$ the field strength of $A$ is $F(a \wedge b) = d\Omega + \Omega \times \Omega + (\text{derivatives of } \hb^{-1}, S, \beta) + \Theta(a) \times \Theta(b)$ plus cross terms $\Omega \times \Theta$; the derivative and cross terms lie in the $\nInf$, $\nO$ and dilation blocks (the torsion, the derivative of the slot and the derivative of $\beta$), because the Lorentz algebra preserves each of those blocks. So the Lorentz block of $F$ is $-\mathcal{R}(a \wedge b)$ (the store's sign, Theorem \claimref{synthesis.gtg-tractor-gauge-field-reproduces-tractor-connection}(iv)) plus the Lorentz block of $\Theta(a) \times \Theta(b)$.

\emph{(i).} Translations commute: $(\nInf \wedge a) \times (\nInf \wedge b) = 0$ and $(\nO \wedge c) \times (\nO \wedge d) = 0$ since $\nInf$, $\nO$ are null and orthogonal to spacetime vectors. The cross terms give $T_{S}(a \wedge b) = a \wedge S(b) + S(a) \wedge b$ in the Lorentz block, using $\nO \cdot \nInf = 1$; the check computes $T_{S}$ on the sixteen basis slots and confirms $T_{\rho\,\id} = 2\rho\, a \wedge b$, which with $\rho = \Lambda/6$ is the stored $\Theta(A) \times \Theta(B) = \tfrac{\Lambda}{3}A \wedge B$ (Theorem \claimref{synthesis.gtg-tractor-gauge-field-reproduces-tractor-connection}(v)). The dilation $\nInf \wedge \nO$ commutes with the Lorentz algebra and maps $\nInf \wedge a$ and $\nO \wedge c$ into their own blocks, so $\beta$ never reaches the Lorentz block. Since $I \wedge V = I \wedge X = \nInf \wedge \nO$ (Theorem \claimref{dark-energy.frozen-point-makes-stelle-west-vector}), $\epsilon(I, V, B, C)$ for Lorentz bivectors $B$, $C$ is the coefficient of $i_{6}$ in $B \wedge C \wedge \nInf \wedge \nO$; for any bivector with a component outside $\Lambda^{2}W$ the wedge with $\nInf \wedge \nO$ either vanishes or has no $i_{6}$ component paired with a Lorentz partner, and the check confirms that the Gram matrix on $\Lambda^{2}W$ is the same at $\Lambda = \tfrac97, 3, -3$. Preservation: $\Theta(a) \times (\nInf + \rho\nO) = \rho\,a - S(a) + \beta(a)(\nInf - \rho\,\nO)$, so it vanishes for all $a$ if and only if $S = \rho\,\id$ and $\beta = 0$; the check solves the linear system for $(S, \beta)$ (twenty unknowns, kernel zero) and tests $\rho\,\id$ and three perturbations. This is the bijection of Theorem \claimref{synthesis.frozen-point-field-space-is-the-gtg-pair-and-its-omega-equation-is-the-torsion-equation}(ii) (connections preserving $\nInf + \rho\nO$ are exactly $\Omega - (\nInf - \rho\nO) \wedge \hb^{-1}$), recomputed here.

\emph{(ii).} With $F_{W} = -\mathcal{R} + T_{S}$, $\mathcal{L} = \epsilon(F_{W}, F_{W})$ is quadratic in $S$: $\mathcal{L} = \epsilon(\mathcal{R}, \mathcal{R}) - 2\epsilon(\mathcal{R}, T_{S}) + \epsilon(T_{S}, T_{S})$. In index reading $\epsilon(T_{S}, T_{S})$ is a multiple of the second elementary symmetric function $\tfrac12((\tr S)^{2} - \tr(S^{2}))$ of $S$, whose polarisation $(S, S') \mapsto \tr S\,\tr S' - \tr(SS')$ is nondegenerate on all linear maps (a kernel element would satisfy $S = (\tr S)\id$, so $\tr S = 4\tr S$ and $S = 0$); the check finds rank $16$. The slot equation is therefore linear with a unique solution. On an algebraic curvature tensor $\mathcal{R}(a \wedge b) = \mathcal{W}(a \wedge b) + \Sch(a) \wedge b + a \wedge \Sch(b)$ (the decomposition of Fact \claimref{synthesis.imported-gtg-dictionary}), so $S = \Sch$ gives $F_{W} = -\mathcal{W}$; since the Weyl part has zero trace on every pair of slots, $\epsilon(\mathcal{W}, T_{S'}) = 0$ for every $S'$ (the $\epsilon$-contraction with $a \wedge S'(b)$ takes a double trace of $\mathcal{W}$), so $S = \Sch$ solves the slot equation and by uniqueness it is the solution. The check confirms $S^{*} = \Sch$ and $F_{W} = -\mathcal{W}$ exactly at two sample algebraic curvatures and at de Sitter curvature ($\mathcal{R}(a \wedge b) = \kappa\, a \wedge b$, $S^{*} = \tfrac{\kappa}{2}\id$). Then $\mathcal{L} = \epsilon(\mathcal{W}, \mathcal{W}) = \tfrac14|W|^{2}$ by Theorem \claimref{synthesis.frozen-point-integrand-on-tractor-transport-is-weyl-squared}(iii), which the check recomputes as a quadratic form on the twenty-dimensional space of algebraic curvature tensors. Since the integrand does not see $\beta$ (part (i)), the dilation may be varied or set to zero; this all-varied field space, restricted to curvature of $\Omega$ obeying the first Bianchi identity, is that of torsion-free conformal Cartan connections (values in all bivectors, nondegenerate soldering, gauged by the parabolic group of the line of $X$). The $\Omega$-equation of the all-varied row with torsion is not computed; with torsion the antisymmetric part of the slot responds (the check finds a nonzero antisymmetric solution at a curvature two-form violating the Bianchi identity), so the statement is restricted to the torsion-free class. Elimination of an auxiliary field whose equation is algebraic and uniquely solvable gives an action for the remaining fields with the same solutions (substitute the solution; the variation with respect to the eliminated field vanishes there); this is the standard argument and is not checked by the script.

\emph{(iii).} Write the slot as $S = S_{0} + \sum_{k} x_{k}P_{k}$ with $S_{0} = \rho\,\id$ when the trace is frozen and $S_{0} = 0$ otherwise, and $P_{k}$ a basis of the varied parts (trace $1$, antisymmetric $6$, trace-free symmetric $9$, all with respect to the metric $\mathrm{diag}(1,-1,-1,-1)$). Then $\mathcal{L} = c(\mathcal{R}) + b(\mathcal{R})^{\top}x + x^{\top}Qx$ with $Q$ the restriction of the nondegenerate form of (ii); on each subset of parts $Q$ is invertible (the check inverts it), and $\mathcal{L}_{\mathrm{red}} = c - \tfrac14 b^{\top}Q^{-1}b$. This is a polynomial of degree two in the twenty components of an algebraic curvature tensor; the check splits it into quadratic, linear and constant parts, fits the quadratic part exactly in the basis $(\tfrac14|W|^{2}$, the Weyl-dual pairing, $|\Ric_{0}|^{2}$, $R^{2})$ of Lorentz-invariant quadratic forms (the fit is exact, with zero coefficient on the dual pairing in every row), and checks that the linear part is a multiple of $R$. With $\tfrac14 E_{4} = \tfrac14|W|^{2} - \tfrac12|\Ric_{0}|^{2} + \tfrac{1}{24}R^{2}$ (the check finds that $\epsilon(\mathcal{R}, \mathcal{R})$ itself has exactly these coefficients), the four rows are as stated, at $\rho = \tfrac{5}{11}$ and $\rho = -\tfrac23$, and at $\rho = 0$ only $\tfrac14 E_{4}$ remains. The ratio of constant to linear coefficient is $-2\Lambda$ with $\Lambda = 6\rho$, so the frozen rows read $\tfrac14 E_{4} - 2\rho(R - 2\Lambda)$; on de Sitter curvature ($R = 4\Lambda$, $E_{4} = R^{2}/6$) this vanishes, as $F_{W} = 0$ requires. Torsion: at a curvature two-form violating the first Bianchi identity the check finds a solution of the slot equation with nonzero antisymmetric part, so the inertness of the antisymmetric part holds only without torsion.

\emph{(iv).} The frozen rows are $\tfrac14 E_{4}$ plus $-2\rho(R - 2\Lambda)$; the Euler density has identically vanishing Euler-Lagrange expression (Fact \claimref{synthesis.imported-lovelock-and-anderson-duchamp-natural-lagrangians}(b)), and $-2\rho(R - 2\Lambda)$ is the Einstein-Hilbert density with cosmological constant $\Lambda$ for $\rho \neq 0$, whose equations are Einstein's; at $\rho = 0$ nothing remains. In the other rows write $\mathcal{L}_{\mathrm{red}} = \tfrac14 E_{4} + L_{1} + q$ with $L_{1}$ the Einstein-Hilbert and volume part (possibly zero) and $q$ the remainder $-\tfrac{1}{24}R^{2}$, $\tfrac12|\Ric_{0}|^{2}$ or $\tfrac12|\Ric_{0}|^{2} - \tfrac{1}{24}R^{2}$. Suppose the Euler-Lagrange expression of $\mathcal{L}_{\mathrm{red}}$ were of second order. Then so is that of $q$, which is natural and divergence-free, so by Fact \claimref{synthesis.imported-lovelock-and-anderson-duchamp-natural-lagrangians}(a) it equals $c_{1}\mathcal{G}(a) + c_{2}a$. Under a constant rescaling of the metric by $\lambda^{2}$ the density $q\,\mathrm{vol}$ is invariant in four dimensions, so its Euler-Lagrange expression scales as $\lambda^{-4}$ (as a map), while $\mathcal{G}$ scales as $\lambda^{-2}$ and $a$ as $1$; hence $c_{1} = c_{2} = 0$ and $q$ is variationally trivial, so by (b) $q = c_{E}E_{4} + c_{P}P_{4}$ plus a divergence, and its integral over a closed four-manifold depends only on the topology. On $S^{2} \times S^{2}$ with product metric of curvatures $K_{1}$, $K_{2}$ (Riemannian reading, as in Theorem \claimref{synthesis.frozen-point-integrand-on-tractor-transport-is-weyl-squared}(v)) the check computes, in units of $16\pi^{2}$: $\int E_{4} = 8$ for every ratio, $\int |\Ric_{0}|^{2} = (K_{1} - K_{2})^{2}/(K_{1}K_{2})$ ($0$, $\tfrac12$, $\tfrac{9}{10}$ at $(1,1)$, $(1,2)$, $(2,5)$), $\int R^{2} = 4(K_{1} + K_{2})^{2}/(K_{1}K_{2})$ ($16$, $18$) and $\int |W|^{2} = \tfrac43(K_{1}+K_{2})^{2}/(K_{1}K_{2})$; each remainder takes different values across the family, a contradiction. So only the frozen rows have second-order equations, and in the all-varied row the equations are Bach's (Fact \claimref{synthesis.imported-conformal-gravity-bach-equations-and-conformal-gauge-form}(a)). The rows with $\beta$ or the antisymmetric part varied give Einstein's equations (on torsion-free fields) but by (i) their connections preserve no tractor for generic values of those fields (with $S = \rho\,\id$ and $\beta \neq 0$ a tractor is parallel only in non-generic cases, where $\beta$ is tied to the frame by a parallel vector); on shell the antisymmetric part vanishes; $\beta$ stays undetermined (a gauge direction of the integrand).

\emph{(v).} The slot one-form $S(\hb^{-1}(a))$ is a field like $\hb^{-1}$ and $\Omega$, and $T_{S}$ is the wedge of the coframe with it, so $\mathcal{L}$ in every row is a wedge polynomial in $\Omega$, the coframe and the slot one-form, without inverse frame and without Hodge dual; it is invariant under the Lorentz rotations and the dilation (which rescales the coframe and the slot inversely and fixes $\nInf \wedge \nO$). In the all-varied row the solution $S = \Sch$ requires the Ricci trace, an inverse frame, but only as the outcome of the slot equation. Clause (vi) of Theorem \claimref{synthesis.curvature-only-tractor-actions-are-macdowell-mansouri} classifies slots built from $(\hb, \Omega)$ (the Lorentz commutant on $W$ is one-dimensional, so $S = \rho\,\id$); its slot step is argued there and proved as clause (i) of Theorem \claimref{synthesis.inverse-frame-free-tractor-gauge-field-preserves-a-frozen-point}, under the premise (T) that the slots are built from $(\hb, \Omega)$ only, so that the slot is not an independent field. That premise, not (P), is what separates the $I$-preserving field space from the unreduced one; with (T), (P) and covariance that forcing gives $S = \rho\,\id$ with $\rho$ constant. In gauge-invariant form, under the parabolic group that gauges the tractor field, (T) with (P) is the statement that the varied connection admits a parallel tractor $I$ with $\ip{X}{I} = 1$ (Premise \claimref{synthesis.field-space-premise-parallel-frozen-point}; by Fact \claimref{synthesis.imported-holonomy-reductions-of-cartan-geometries}(a) a parallel tractor is a holonomy reduction of the connection to the stabiliser of $I$): a special-conformal rotor moves the frozen-row connection to one whose slot is not $\rho\,\id$ and which preserves the rotated tractor instead, so $S = \rho\,\id$, $\beta = 0$ is that statement in the gauge where $I$ is constant. The Einstein condition is not presupposed anywhere: the frozen value is independent of the curvature, and the frozen and normal rows meet where $\Sch = \rho\,\id$, the Einstein fields with $\Lambda = 6\rho$ (Theorem \claimref{synthesis.gtg-tractor-gauge-field-reproduces-tractor-connection}(ii)).
